% ================================================================
%  conteudo/referencias.tex
%  ABNT NBR 6023:2018
%
%  Regras:
%    - Alinhadas à esquerda, sem recuo
%    - Espaçamento simples dentro de cada referência
%    - Uma linha em branco entre referências
%    - Título da obra em negrito
%    - Ordenadas alfabeticamente pelo sobrenome do primeiro autor
% ================================================================

\clearpage
\addcontentsline{toc}{section}{Referências Bibliográficas}
\vspace{1\baselineskip}
{\centering\bfseries REFERÊNCIAS BIBLIOGRÁFICAS\par}
\begingroup
\vspace{-0.5\baselineskip}
\singlespacing
\setlength{\parskip}{\baselineskip}
\setlength{\parindent}{0pt}

BERES, C. \textbf{Aproveitamento de resíduos da indústria vinícola}. Rio de Janeiro: UFRJ, 2013.

CÂMARA, M. C. \textbf{Técnicas de Microbiologia Industrial}. São Paulo: Editora Científica, 2023.

CARVALHO, W. \textbf{Biocatalisadores e processos fermentativos}. Lorena: USP, 2010.

ESCOLA ESTADUAL DE EDUCAÇÃO PROFISSIONAL (EEEP). \textbf{Microbiologia de Alimentos}. Fortaleza: SEDUC, 2012.

FARINAS, C. S.; BARBOZA, M. \textbf{Fungos filamentosos na biotecnologia}. São Carlos: Embrapa, 2012.

SOARES, L.; ISHIMOTO, E. Y.; BATTESTIN, V. \textbf{Isolamento e seleção de microrganismos}. Campinas: UNICAMP, 2013.

VIEIRA, G.; FERNANDES, P. \textbf{Bioquímica de Leveduras}. Lisboa: IST, 2012.
\endgroup